\documentclass[../main]{subfiles}
\begin{document}
\chapter*{Preface}
This lecture note is written to address the absence of a well-formatted, illustrated reference for this course.  
It began as spontaneous notes taken during class, and has gradually evolved through the careful revisions of the instructor and teaching assistants, as well as polishing assisted by large language models.

Some parts may remain incomplete or insufficiently elaborated. I sincerely ask for your understanding.  
If you find any typos, inaccuracies, or suggestions for improvement, you are warmly welcome to open an issue on GitHub:

\begin{center}
\url{https://github.com/PhotonYan/MachineLearningCourse2025}
\end{center}

\begin{flushright}
    \textsc{Shaoheng Yan} \\[3pt]
    \scriptsize\textit{Tong Class 2024, Yuanpei College} \\
    \textit{M$\mu$ Lab, Institute for Artificial Intelligence (IAI)} \\[3pt]
    \textit{Peking University, November 2025}
\end{flushright}
\begin{center}
    { \large \textbf{Environments in This Note}}
\end{center}


This lecture note adopts a unified visual style for mathematical structures and remarks.  
Below are the environments used throughout the text.

\begin{definition}
\end{definition}

\begin{theorem}
\end{theorem}

\begin{proposition}
\end{proposition}

\begin{remark}
    Remark represents supplementary comments made by the instructor during class.
\end{remark}

\begin{note}
    Note indicates additional reflections by the author, which the reader may skip if desired.
\end{note}

\begin{example}
\end{example}

% \begin{exercise}
% \end{exercise}

\begin{proof}
\end{proof}
\begin{solution}
    
\end{solution}

\end{document}